\section*{Exercises}
Exercises~\ref{ex:6.1}--\ref{ex:6.4} are core exercises, and Exercises~\ref{ex:6.5} and~\ref{ex:6.6} are transfer exercises. Complete solutions are in Appendix~\ref{app:sola}, under ``Week~6.''
\begin{exercise}\label{ex:6.1}List the members of $[-2]_5$ that lie between $-12$ and $12$ inclusive. Which integer in $\{0,1,2,3,4\}$ represents the same class?\end{exercise}
\begin{exercise}\label{ex:6.2}Prove that multiplication of residue classes, defined by $[a]_m\cdot[b]_m=[ab]_m$, is well-defined. That is, show that if $[a']_m=[a]_m$ and $[b']_m=[b]_m$, then $[a'b']_m=[ab]_m$.\end{exercise}
\begin{exercise}\label{ex:6.3}For the four seats around the table in Section~\ref{ch06:sec:structure}, show that the rotation $r_2$ by two seats is its own inverse. Explain why this does not make $r_2$ the identity.\end{exercise}
\begin{exercise}\label{ex:6.4}Is ``has the same number of elements as'' an equivalence relation on a collection of finite sets? Check each of the three properties.\end{exercise}
\begin{exercise}\label{ex:6.5}Let $m$ and $d$ be positive integers. Prove that the rule $[a]_m\mapsto[a]_d$ is a well-defined function from $\ZZ/m\ZZ$ to $\ZZ/d\ZZ$ exactly when $d$ divides $m$.\end{exercise}
\begin{exercise}\label{ex:6.6}Give a family of three sets, ordered by inclusion, that has exactly two maximal members and no maximum. Then add one set to the family so that the new family has a maximum.\end{exercise}
